% S8a threat model and definitions. Must stay consistent with
% docs/research-contract.md (revision 3) §§4, 5, 8. The IEEE S&P 2027 version
% of this section is at commit 96a8c7c.
\section{Threat Model and Definitions}\label{sec:threat}

\para{System.}
A ZK application has three artifacts. The \emph{circuit} defines a relation
$R_v(p,w)$ between a public statement $p$ and a witness $w$ for verification
configuration $v$ (circuit, compiled constraints and key). The
\emph{generated verifier} accepts a proof for a public vector $x$; the actual
decoding relation $D_v(p,x)$ fixes the arity, order and encoding of $x$. The
\emph{consumer} is the contract that assembles $x$, calls the verifier and
commits effects. We write $T(s,d,\tau,s')$ for a successful consumer
execution from state $s$ to $s'$ with calldata $d$ and committed trace
$\tau$.%
% src: docs/research-contract.md#4-system-model
\
\TODO{S0: contract revision 4 must record the Circom/Groth16 scope and the
Noir/Honk generator survey; revision 3 still names a Noir/ACIR producer.}

\para{Adversary.}
The adversary may choose witnesses (and deviate from the honest witness
generator), public inputs, proof bytes, calldata and transaction order; call
every public entry point and role that the evaluated deployment allows; reuse
any proof or value it observes unless the protocol prevents it; and start from
any reachable pre-state. Capabilities that a case needs beyond these are
recorded per case. For example, the historical Payy defect requires a
registered prover~\cite{payy-ghsa}, while Aztec Connect's last contract
version let any caller submit a rollup.%
% src: docs/research-contract.md#51-adversary-capabilities; reports/AZ1/REPORT.md (RollupProcessorV3.sol:665-675)

\para{Trust assumptions.}
We assume that the proof system is sound for the pinned relation and key;
that the analyzed source corresponds to the compiled constraints, the key and
the consumer configuration; that stated hash functions are collision
resistant; that the reviewed policy is correct; and that token contracts
follow atomic ERC-20/721 semantics. We do \emph{not} assume that a public
value is constrained because it should be, or that an effect is bound because
it looks bound. Out of scope are breaking the proof system, consensus,
compiler or hardware faults, privacy leaks unrelated to execution, liveness,
non-standard tokens, and any interaction with live deployments. Unsupported
constructs (proxies, inline assembly, delegate calls, reentrancy) yield \UNK{}.
A key from an uncontributed setup ($\gamma=\delta$) is a trusted-setup failure
under the research contract; we report it as a verifier-configuration class
(P4) and a root cause, never as evidence that joint analysis is needed.%
% src: docs/research-contract.md#52-trusted-assumptions, #54-initially-out-of-scope

\para{Effects, permissions and obligations.}
A reviewed policy $\sigma$ states which consumer operations are
\emph{protected effects} (transfers, mints, burns, credential or nullifier
writes, state-root updates), which values of an accepted statement are
\emph{permissions}, and which \emph{obligations} an accepted transition
creates (for example, collecting a deposit that the proof credited). Each
record carries a kind, identity, asset, quantity, beneficiary and domain. For
an execution we write $P$ for the multiset of permissions, $O$ for the
required obligations and $E$ for the committed protected effects. The policy
names roles and protected operations only. It never states that a value is
constrained, that an effect is bound, or the expected verdict.%
% src: docs/research-contract.md#4-system-model, #6-non-circular-specification-boundary

\para{Binding.}
An execution is \emph{authorized} if every effect in $E$ is matched by a
distinct compatible permission in $P$, preserving identity, asset,
beneficiary, domain and quantity, and no permission is used beyond its
multiplicity. It is \emph{covered} if every obligation in $O$ is discharged by
a compatible committed effect, validly cancelled, or turned into a pending
obligation that the protocol explicitly permits. A \emph{binding flaw} is a
reachable state $s$, witness $w$ and calldata $d$ such that $R_v(p,w)$ and
$D_v(p,x)$ hold, the unchanged verifier accepts, the unchanged consumer
commits $T(s,d,\tau,s')$, and authorization or coverage fails.%
% src: docs/research-contract.md#41-authorization, #42-coverage, #4-system-model (joint counterexample)
\
The one-sentence form used in \S\ref{sec:intro} is the intuition behind
these definitions. If the effects can change while $p$ and $s$ stay fixed,
some effect parameter came from outside the proof, and one of the two
executions is unauthorized unless the policy explicitly lets that parameter
vary (for example, a faucet that pays any caller). If the same $p$ can be accepted again, the
multiplicity condition fails. Coverage adds what the one-sentence form
leaves implicit: an accepted proof can also \emph{create} obligations that
the consumer must discharge.

\para{Proof coverage and extent.}
For a sequence $x[0..L)$ fixed by the public inputs (directly or through a
commitment in them), the proof \emph{covers} index $i$ if $R_v$ constrains
$x[i]$. Extent binding asks whether the range the consumer acts on equals the
covered range. It fails in direction~(a), \emph{reliance beyond binding}, when
the consumer relies on some $x[i]$ that the proof does not cover. It fails in
direction~(b), \emph{binding beyond discharge}, when the proof covers and
applies a non-neutral $x[i]$ for which the consumer discharges no
obligation, because the consumer's bound comes from a value the proof does not
fix. Direction~(a) breaks authorization; direction~(b) breaks coverage.%
% src: reports/AZ1/REPORT.md#extent-binding-pattern-definition

\para{Verdicts.}
Every analysis returns one of three verdicts. \VIOL{}: a model that satisfies
the supported assumptions violates authorization or coverage. \VWS{}: no
violation exists within the declared model, bounds and assumptions, and the
check is non-vacuous. \UNK{}: extraction, modeling, provenance, the solver or
validation is insufficient. A timeout, a build failure, an unsupported
feature or a missing artifact is \UNK{} and never counts as safe. Each
verdict carries a reason code, the bounds and the assumptions it rests on;
an unresolved specification is an \UNK{} reason, not a fourth verdict.%
% src: docs/research-contract.md#8-required-result-categories

\para{Evidence ladder.}
Each candidate also carries one evidence grade. \textbf{E0}: no report.
\textbf{E1}: a tool warning. \textbf{E2}: a model that satisfies the actual
compiled relation, the decoder and the consumer model. \textbf{E3}: a real
proof of that assignment, accepted by the pinned, unchanged verifier.
\textbf{E4}: that proof, accepted by the unchanged consumer from an admitted
or reachable state, produces an unauthorized committed effect or leaves a
required obligation undischarged, checked on state. Only E4 is an
\emph{end-to-end validated} binding flaw. A fixed admitted pre-state is a
narrower mode than all reachable states, and we label it when we use it. The
same validation process, budget and adjudication apply to \tool{} and to every
baseline.%
% src: docs/research-contract.md#8 (evidence ladder), #4 (fixed admitted pre-state)
